\documentclass[10pt,a4paper]{article}
\usepackage[margin=15mm]{geometry}
\usepackage{fontspec}
\usepackage{xeCJK}
\usepackage{booktabs}
\usepackage{multirow}
\usepackage{siunitx}
\usepackage{caption}
\usepackage{pdflscape}
\usepackage{graphicx}
\usepackage{array}
\usepackage{xurl}
\usepackage[hidelinks]{hyperref}

\setmainfont{Times New Roman}
\setsansfont{Arial}
\setCJKmainfont{Microsoft YaHei}
\setCJKmonofont{Microsoft YaHei}
\sisetup{detect-weight=true,table-number-alignment=center}
\captionsetup{font=small,labelfont=bf,skip=4pt}
\setlength{\parindent}{0pt}
\setlength{\parskip}{0.45em}
\setlength{\tabcolsep}{4pt}
\pagestyle{plain}
\sloppy

\newcommand{\uprep}{$R_{\mathrm{up}}^{\mathrm{rep}}$}
\newcommand{\upwho}{$R_{\mathrm{up}}^{\mathrm{who}}$}
\newcommand{\onwho}{$R_{\mathrm{on}}^{\mathrm{who}}$}
\newcommand{\na}{---}
\newcommand{\zh}[1]{\par\noindent{\small #1}\par}

\begin{document}

\begin{center}
{\LARGE\bfseries Experiments / 实验}\\[4pt]
{\large PHAST-DP Main Results, Ablation Studies, and Baseline Comparisons}\\[2pt]
{\large PHAST-DP 主模型结果、消融实验与基线对比}
\end{center}

\section{Research Questions / 研究问题}
The experiments are designed to answer the following four research questions.
\zh{本文实验围绕以下四个研究问题展开。}

\begin{enumerate}
\item \textbf{RQ1:} Can PHAST-DP accurately identify station-level bottleneck
events at one-minute resolution across Start$\leq5$, Start$\leq10$, and
Start$\leq15$ prediction ranges?
\zh{\textbf{RQ1：}PHAST-DP 能否以一分钟分辨率，在 Start$\leq5$、Start$\leq10$ 和 Start$\leq15$ 三种预测范围内准确识别工位级瓶颈事件？}

\item \textbf{RQ2:} How accurately can the model estimate event onset,
duration, remaining order time, and process cause while distinguishing ongoing
from upcoming bottlenecks?
\zh{\textbf{RQ2：}在区分进行中与即将发生瓶颈的同时，模型对事件起点、持续时间、订单剩余时间和过程原因的估计有多准确？}

\item \textbf{RQ3:} Do graph structure and station-group embeddings contribute
to station identification, early-warning recall, remaining-time estimation, and
cause classification?
\zh{\textbf{RQ3：}图结构和工位分组嵌入是否有助于工位识别、提前预警召回、剩余时间估计和原因分类？}

\item \textbf{RQ4:} What performance is obtained when clustering is disabled
and mixed-disturbance episodes are introduced, and can the effects of these two
changes be separated?
\zh{\textbf{RQ4：}关闭聚类并引入混扰 episode 后能够取得怎样的性能？数据混合和去除聚类这两项变化的影响能否被分别识别？}
\end{enumerate}

\section{Data and Protocol / 数据与实验协议}

\subsection{Task and Evaluation Protocol / 任务与评测协议}
The experiments evaluate station-level bottleneck-event prediction at three start
horizons, Start$\leq 5$, Start$\leq 10$, and Start$\leq 15$ min, while requiring
each target event to last at least eight minutes (Min8). The model observes a
30-min history and evaluates a 20-grid future label range at a one-minute
resolution. An \emph{ongoing} target is already active at the last history grid;
an \emph{upcoming} target starts after the observation boundary.
\zh{实验在三个起始时间范围内评估工位级瓶颈事件预测：Start$\leq5$、Start$\leq10$ 和 Start$\leq15$ 分钟，同时要求每个目标事件至少持续 8 分钟（Min8）。模型观察 30 分钟历史，并以 1 分钟为时间分辨率，在未来 20 个网格的标签范围内进行评测。若目标在历史序列最后一个网格已经发生，则记为进行中（ongoing）；若目标在观测边界之后才开始，则记为即将发生（upcoming）。}

The station-identification metrics are \texttt{will15}/\texttt{who} precision,
recall, and F1. The name \texttt{will15} is retained by the implementation but
refers to the active Start cap, rather than always denoting a fixed 15-min task.
The stricter report metric counts a true positive only when the station is
correct and $|\hat{s}-s|\leq 3$ min. Thus, \upwho measures station detection,
whereas \uprep additionally enforces start-time tolerance. Remaining-order-time
and event-duration errors are reported as mean absolute error (MAE; lower is
better). Process-cause performance is reported by accuracy, macro recall, and
per-class recall.
\zh{工位识别指标包括 \texttt{will15}/\texttt{who} 的精确率、召回率和 F1。代码沿用了 \texttt{will15} 这一名称，但该指标实际对应当前采用的 Start 上限，并不总是表示固定的 15 分钟任务。更严格的 report 指标只有在工位预测正确且 $|\hat{s}-s|\leq3$ 分钟时才计为真阳性。因此，\upwho 衡量是否识别出正确工位，而 \uprep 还额外要求起始时间满足容差。订单剩余时间和事件持续时间采用平均绝对误差（MAE，越低越好）；过程原因识别则报告准确率、宏平均召回率和逐类别召回率。}

All model and threshold decisions use validation data. The selected checkpoint
and threshold are then frozen for test evaluation. Results from different Start
caps are not pooled because their positive sets and upcoming denominators differ.
All reported runs use seed 42; consequently, no claim of statistical
significance is made.
\zh{所有模型选择和阈值选择均只使用验证集。确定 checkpoint 和阈值后，将二者冻结并用于测试集评测。由于不同 Start 上限对应的正例集合和 upcoming 分母不同，因此不合并不同 Start 档位的结果。所有实验均只使用 seed 42，所以本文不声称结果具有统计显著性。}

\subsection{Data Splits and Experiment Families / 数据划分与实验系列}
\begin{center}
\captionof{table}{Experiment families and protocol boundaries; episode counts are train/validation/test. / 实验系列及协议边界；episode 数依次为训练/验证/测试。}
\small
\resizebox{\textwidth}{!}{
\begin{tabular}{llllll}
\toprule
Family & Dataset & Episodes & Main selection metric & Training budget & Test status \\
\midrule
Final PHAST-DP (\texttt{12\_3})
 & \texttt{dense\_i1} & 140/28/36
 & validation \texttt{will15\_f1} & 100 epochs, patience 40 & evaluated \\
Earlier PHAST-DP (\texttt{12\_2})
 & \texttt{dense\_i1} & 140/28/36
 & validation \texttt{report\_f1} & 100 epochs & evaluated \\
No-clustering mixed-data experiment
 & \texttt{dense\_i1\_mix} & 166/37/45
 & validation \texttt{will15\_f1} & 100 epochs, patience 40 & evaluated \\
Structure ablations
 & \texttt{dense\_i1} & 140/28/36
 & validation \texttt{will15\_f1} & 50 epochs, from scratch & evaluated \\
B4 BTGCN / B5 BSTAN
 & fixed 208-episode set & 138/30/40
 & validation \texttt{will15\_f1} & 100 epochs, patience 40 & test unused \\
\bottomrule
\end{tabular}}
\end{center}

\section{Implementation Details / 实现细节}
The PHAST-DP encoder contains five blocks with an embedding dimension of 80
and an eight-dimensional Laplacian positional encoding. Graph, semantic, and
temporal attention use 2, 2, and 4 heads, respectively; these are separate
attention modules and should not be described as a single eight-head module.
The event head has hidden dimension 80. Training uses batch size 16, a 30-min
input history, one-minute windows, and a 20-grid future label range.
\zh{PHAST-DP 编码器包含 5 个块，嵌入维度为 80，并使用 8 维拉普拉斯位置编码。图注意力、语义注意力和时间注意力分别使用 2、2 和 4 个头；它们属于不同的注意力模块，因此不能表述为一个统一的“8 头注意力模块”。事件预测头的隐藏维度为 80。训练 batch size 为 16，输入历史长度为 30 分钟，时间窗口为 1 分钟，未来标签范围为 20 个网格。}

Optimization uses AdamW with a fine-tuning learning rate of $5\times10^{-5}$,
weight decay $0.05$, cosine learning-rate decay, and minimum learning rate
$10^{-6}$. The final runs allow at most 100 epochs and use early-stopping
patience 40. Event-positive windows are oversampled by a factor of 4. The model
is initialized from the earlier recall-oriented
\texttt{dense\_i1\_a1\_prefix8} checkpoint, unfreezes the last five encoder
blocks and occupancy decoder, and jointly optimizes event, occupancy,
remaining-time, and cause objectives. STGNPP and the legacy
\texttt{score}/\texttt{will}/\texttt{mark}/\texttt{tts} heads are disabled in
the retained configuration.
\zh{优化器采用 AdamW，微调学习率为 $5\times10^{-5}$，权重衰减为 $0.05$，使用余弦学习率衰减，最低学习率为 $10^{-6}$。最终实验最多训练 100 个 epoch，早停 patience 为 40；事件正例窗口按 4 倍过采样。模型从较早的召回导向 checkpoint \texttt{dense\_i1\_a1\_prefix8} 初始化，解冻最后 5 个编码器块和占用解码器，并联合优化事件、占用、剩余时间和原因任务。最终配置关闭 STGNPP，以及旧版 \texttt{score}/\texttt{will}/\texttt{mark}/\texttt{tts} 预测头。}

The final PHAST-DP configurations are
\texttt{FactoryBN\_dense\_12\_3\_start\{5,10,15\}\_min8\_opt.json}.
They use the shared threshold set
$\{0.45,0.50,0.55,0.60,0.65,0.70,0.75,0.80,0.85,0.90,0.94\}$,
the validation gate $P\geq0.80$ and $R\geq0.70$, Start-specific upcoming/far-start
positive weights of 9/10/11, hard false-positive weighting, occupancy union, and
progress-weighted remaining-time training. Start10 extends Start5 and Start15
extends Start10. The prediction threshold is selected independently for each
checkpoint but never on test data.
\zh{最终 PHAST-DP 使用配置 \texttt{FactoryBN\_dense\_12\_3\_start\{5,10,15\}\_min8\_opt.json}。三档实验共享阈值候选集合 $\{0.45,0.50,0.55,0.60,0.65,0.70,0.75,0.80,0.85,0.90,0.94\}$，验证集门槛为 $P\geq0.80$ 且 $R\geq0.70$；Start5/10/15 的 upcoming 与远起点正例权重分别为 9/10/11，并使用困难假阳性加权、占用并集和按生产进度加权的剩余时间训练。Start10 继承 Start5，Start15 再继承 Start10。每个 checkpoint 独立选择预测阈值，但测试集从不参与该选择。}

\section{Main Results / 主实验结果}
\subsection{Selection of the Retained Main-Model Family / 最终主模型系列的选择}
The earlier \texttt{12\_2} runs and the later protocol-aligned \texttt{12\_3}
runs share the same \texttt{dense\_i1} split and Min8/Start task. The main
difference is that \texttt{12\_3} switches checkpointing to
\texttt{will15\_f1}, applies the common precision/recall gate and threshold grid,
increases far-upcoming supervision, and introduces the progress-weighted
remaining-time objective. Table~\ref{tab:model-selection} records the complete
selection evidence.
\zh{较早的 \texttt{12\_2} 与后续完成协议对齐的 \texttt{12\_3} 使用相同的 \texttt{dense\_i1} 数据划分和 Min8/Start 任务。主要区别是：\texttt{12\_3} 改用 \texttt{will15\_f1} 选模，采用统一的精确率/召回率门槛和阈值集合，加强远距离 upcoming 监督，并加入按生产进度加权的剩余时间目标。表~\ref{tab:model-selection} 给出了完整的选择依据。}

\begin{center}
\captionof{table}{Earlier \texttt{12\_2} versus protocol-aligned \texttt{12\_3}; selection is based on validation behavior and protocol consistency, not one test cell. / 较早的 \texttt{12\_2} 与协议对齐的 \texttt{12\_3}；选择依据是验证表现和协议一致性，而非单个测试数值。}
\label{tab:model-selection}
\scriptsize
\resizebox{\textwidth}{!}{
\begin{tabular}{llrrrrrrrrrr}
\toprule
Start & Family & Val $P$ & Val $R$ & Val F1 & Test $P$ & Test $R$ &
Test F1$_{\mathrm{who}}$ & Test F1$_{\mathrm{rep}}$ &
\uprep\ (test) & \onwho\ (test) & Best ep. \\
\midrule
$\leq5$  & 12\_2 & .865 & .872 & .869 & .832 & .808 & .820 & .819 & .542 & .923 & 33 \\
$\leq5$  & 12\_3 & .837 & .891 & .863 & .794 & .817 & .805 & .804 & .539 & .937 & 39 \\
\midrule
$\leq10$ & 12\_2 & .805 & .836 & .820 & .750 & .763 & .757 & .751 & .550 & .923 & 83 \\
$\leq10$ & 12\_3 & .831 & .840 & .835 & .806 & .761 & .783 & .778 & .528 & .936 & 41 \\
\midrule
$\leq15$ & 12\_2 & .814 & .803 & .808 & .772 & .716 & .743 & .741 & .489 & .925 & 37 \\
$\leq15$ & 12\_3 & .841 & .835 & .838 & .806 & .742 & .772 & .768 & .534 & .925 & 42 \\
\bottomrule
\end{tabular}}
\end{center}

\texttt{12\_2} is stronger at Start$\leq5$, but \texttt{12\_3} is stronger at
Start$\leq10$ and Start$\leq15$, has higher mean validation F1
($0.845$ versus $0.832$), and follows the protocol used by B4/B5. Therefore,
\texttt{12\_3} is retained as the final main-model family. This decision avoids
cherry-picking a different recipe from test performance at each horizon.
\zh{\texttt{12\_2} 在 Start$\leq5$ 时更强，但 \texttt{12\_3} 在 Start$\leq10$ 和 Start$\leq15$ 时更优，其平均验证 F1 也更高（$0.845$ 对 $0.832$），并且与 B4/B5 使用的协议一致。因此，最终保留 \texttt{12\_3} 作为主模型系列。该决定避免根据各档测试集表现分别挑选不同配方而造成测试集选择偏差。}

\begin{center}
\captionof{table}{All recorded auxiliary test metrics for the non-retained \texttt{12\_2} family. / 未入选的 \texttt{12\_2} 系列全部已记录辅助测试指标。}
\scriptsize
\resizebox{\textwidth}{!}{
\begin{tabular}{lrrrrrrrrrrrr}
\toprule
Start & Remain MAE & Dur. MAE & Start MAE & F1@1 & F1@2 & F1@3 &
Cause acc. & Cause macro & \upwho & Threshold & $n_{\rm up}/n_{\rm on}/n_{\rm all}$ \\
\midrule
$\leq5$  & 10.64 & 1.93 & .13 & .801 & .814 & .819 & .965 & .969 & .545 & .90 & 336/766/1102 \\
$\leq10$ &  9.61 & 1.77 & .24 & .730 & .749 & .751 & .965 & .970 & .560 & .88 & 604/768/1372 \\
$\leq15$ & 10.20 & 1.88 & .28 & .706 & .735 & .741 & .969 & .974 & .489 & .94 & 708/770/1478 \\
\bottomrule
\end{tabular}}
\end{center}

For completeness, the selected \texttt{12\_2} validation
remain-MAE/cause-accuracy/cause-macro triplets were
$12.03/.952/.769$, $10.86/.956/.774$, and $11.77/.956/.776$ for
Start$\leq5/10/15$. Test cause support was 579 and the train-derived
majority-class accuracy was $0.409$. All three runs completed 100 epochs.
\zh{为保证数据完整性，入选的 \texttt{12\_2} checkpoint 在 Start$\leq5/10/15$ 下的验证集“剩余时间 MAE/原因准确率/原因宏召回率”分别为 $12.03/.952/.769$、$10.86/.956/.774$ 和 $11.77/.956/.776$。测试集原因标签支持数为 579，依据训练集得到的多数类基线准确率为 $0.409$。三组实验均训练满 100 个 epoch。}

\clearpage
\subsection{Final Main-Model Performance / 最终主模型表现}
\subsubsection{Station-Level Bottleneck Prediction / 工位级瓶颈预测}
\begin{center}
\captionof{table}{Final PHAST-DP (\texttt{12\_3}) station-level results; validation is the selected epoch and test uses the frozen checkpoint/threshold. / 最终 PHAST-DP（\texttt{12\_3}）工位级结果；验证列来自入选 epoch，测试列使用冻结的 checkpoint 和阈值。}
\label{tab:main-event}
\scriptsize
\resizebox{\textwidth}{!}{
\begin{tabular}{llrrrrrrrrrr}
\toprule
Start & Split & $P$ & $R$ & F1$_{\mathrm{who}}$ & F1$_{\mathrm{rep}}$ &
\uprep & \onwho & $n_{\mathrm{up}}$ & $n_{\mathrm{on}}$ & $n_{\mathrm{all}}$ & Epoch \\
\midrule
$\leq5$  & Val  & .837 & .891 & .863 & .863 & .739 & .969 & \na & \na & \na & 39 \\
$\leq5$  & Test & .794 & .817 & \bfseries .805 & .804 & .539 & .937 & 336 & 766 & 1102 & 39 \\
\midrule
$\leq10$ & Val  & .831 & .840 & .835 & .835 & .711 & .960 & \na & \na & \na & 41 \\
$\leq10$ & Test & .806 & .761 & \bfseries .783 & .778 & .528 & .936 & 604 & 768 & 1372 & 41 \\
\midrule
$\leq15$ & Val  & .841 & .835 & .838 & .838 & .722 & .960 & \na & \na & \na & 42 \\
$\leq15$ & Test & .806 & .742 & \bfseries .772 & .768 & .534 & .925 & 708 & 770 & 1478 & 42 \\
\bottomrule
\end{tabular}}
\end{center}

The test station-level F1 decreases from $0.805$ to $0.772$ as the task includes
more distant upcoming events. Ongoing recall remains between $0.925$ and $0.937$,
whereas upcoming report recall remains around $0.53$. For Start$\leq15$, the
upcoming bucket recalls are $0.574$ for Start$\leq5$, $0.522$ for starts 6--10,
and $0.433$ for starts above 10 min. This identifies distant early warning,
rather than ongoing detection, as the principal limitation.
\zh{随着任务纳入更远的 upcoming 事件，测试集工位级 F1 从 $0.805$ 下降到 $0.772$。进行中召回率保持在 $0.925$--$0.937$，而 upcoming 严格报告召回率约为 $0.53$。在 Start$\leq15$ 设置下，Start$\leq5$、6--10 分钟以及大于 10 分钟三个 upcoming 桶的召回率分别为 $0.574$、$0.522$ 和 $0.433$。这说明主要限制是远距离提前预警能力，而不是进行中事件检测。}

\subsubsection{Remaining Time and Cause Classification / 剩余时间与原因分类}
\begin{center}
\captionof{table}{Final PHAST-DP auxiliary-head results; remaining-time MAE is in minutes. / 最终 PHAST-DP 辅助任务结果；剩余时间 MAE 的单位为分钟。}
\scriptsize
\begin{tabular}{llrrrr}
\toprule
Start & Split & Remain MAE & Cause accuracy & Cause macro recall & Cause support \\
\midrule
$\leq5$  & Val  & 12.3 & .956 & .963 & \na \\
$\leq5$  & Test & 10.9 & .965 & .967 & 579 \\
$\leq10$ & Val  & 12.6 & .958 & .968 & \na \\
$\leq10$ & Test & 11.2 & .962 & .965 & 579 \\
$\leq15$ & Val  & 11.0 & .967 & .978 & \na \\
$\leq15$ & Test &  9.7 & .974 & .978 & 579 \\
\bottomrule
\end{tabular}
\end{center}

\begin{center}
\captionof{table}{Final PHAST-DP per-class cause recall on test; the final two classes have no test support. / 最终 PHAST-DP 测试集逐类别原因召回率；最后两类在测试集中无样本。}
\scriptsize
\begin{tabular}{lrrrrrr}
\toprule
Start & Transport delay & Material shortage & Starved upstream &
Queue buildup & Blocked downstream & High utilization \\
\midrule
$\leq5$  & .941 & .976 & .966 & .984 & no support & no support \\
$\leq10$ & .926 & .988 & .962 & .984 & no support & no support \\
$\leq15$ & .978 & .988 & .962 & .984 & no support & no support \\
\bottomrule
\end{tabular}
\end{center}

The cause head has four supported classes out of six configured classes.
The majority-class baseline is derived from train only and equals $0.409$.
The unsupported classes are retained in the table to make the support boundary
explicit rather than inflating a six-class claim.
\zh{原因识别头虽然配置了 6 个类别，但只有 4 个类别拥有有效样本。多数类基线只依据训练集确定，其准确率为 $0.409$。表中仍保留两个无测试样本的类别，以明确支持度边界，避免将实际四类有效评测夸大为完整六类评测。}

\clearpage
\section{Ablation Studies / 消融实验}

\subsection{No-Clustering Mixed-Disturbance Experiment / 无聚类混扰数据实验}
This experiment uses
\texttt{FactoryBN\_dense\_mix\_start\{5,10,15\}\_min8\_opt.json} and
\texttt{dense\_i1\_mix}, which contains 248 episodes: the quality-filtered
single-disturbance data plus 55 mixed 2/3/4-dimension episodes. The inherited
configuration has $w_{\mathrm{cluster}}=0$ and does not use a cluster objective.
However, because the dataset also changes, this is a \emph{compound,
ablation-style experiment}, not a pure causal estimate of clustering alone.
The split is 166/37/45, selection uses validation \texttt{will15\_f1}, and the
curriculum is Start5$\rightarrow$Start10$\rightarrow$Start15.
\zh{本实验使用 \texttt{FactoryBN\_dense\_mix\_start\{5,10,15\}\_min8\_opt.json} 和 \texttt{dense\_i1\_mix}。该数据集共含 248 个 episode，由通过质量筛选的单扰动数据和 55 个二维/三维/四维混扰 episode 组成。继承配置中 $w_{\mathrm{cluster}}=0$，不使用聚类目标。但由于该实验同时改变了数据集，所以它属于复合型、类似消融的实验，而不是对“去掉聚类”这一单一因素的严格因果估计。数据划分为 166/37/45，使用验证集 \texttt{will15\_f1} 选模，并按 Start5$\rightarrow$Start10$\rightarrow$Start15 进行课程式训练。}

\begin{center}
\captionof{table}{No-clustering mixed-data event metrics; Start10/15 use the ungated validation best because the gate did not open. / 无聚类混扰数据事件指标；Start10/15 因门槛未开启而采用无门槛验证集最优值。}
\scriptsize
\resizebox{\textwidth}{!}{
\begin{tabular}{llrrrrrrrrrrrr}
\toprule
Start & Split & $P_{\mathrm{rep}}$ & $R_{\mathrm{rep}}$ & F1$_{\mathrm{rep}}$ &
$P_{\mathrm{who}}$ & $R_{\mathrm{who}}$ & F1$_{\mathrm{who}}$ &
\uprep & \upwho & \onwho & Support $n_{\rm up}/n_{\rm on}/n_{\rm all}$ &
Best/total ep. \\
\midrule
$\leq5$  & Val  & .822 & .718 & .766 & .824 & .720 & .769 & .349 & \na & .898 & \na & 1/41 \\
$\leq5$  & Test & .772 & .701 & .735 & .774 & .702 & .737 & .299 & .301 & .888 & 562/1216/1778 & 1/41 \\
\midrule
$\leq10$ & Val  & .728 & .682 & .704 & .742 & .695 & .718 & .460 & \na & .872 & \na & 58/98 \\
$\leq10$ & Test & .684 & .613 & .646 & .697 & .625 & .659 & .324 & .344 & .855 & 1000/1221/2221 & 58/98 \\
\midrule
$\leq15$ & Val  & .746 & .657 & .699 & .760 & .669 & .712 & .445 & \na & .868 & \na & 30/70 \\
$\leq15$ & Test & .697 & .586 & .637 & .704 & .592 & .643 & .320 & .325 & .847 & 1168/1222/2390 & 30/70 \\
\bottomrule
\end{tabular}}
\end{center}

\begin{center}
\captionof{table}{All recorded no-clustering mixed-data auxiliary metrics; time errors are in minutes. / 无聚类混扰数据实验的全部辅助指标；时间误差单位均为分钟。}
\scriptsize
\resizebox{\textwidth}{!}{
\begin{tabular}{lrrrrrrrrrrrrr}
\toprule
Start & Remain & Remain$_w$ & Remain primary & Dur. & Start &
F1@1 & F1@2 & F1@3 & Cause acc. & Cause macro & $n_{\rm cause}$ &
Majority & Threshold \\
\midrule
$\leq5$  & 11.85 & 9.71 & 9.56 & 2.54 & .13 & .712 & .729 & .735 & .944 & .964 & 964 & .495 & .94 \\
$\leq10$ & 10.95 & 9.12 & 9.77 & 2.18 & .27 & .626 & .643 & .646 & .924 & .940 & 964 & .495 & .94 \\
$\leq15$ & 10.92 & 9.03 & 9.57 & 2.26 & .22 & .623 & .636 & .637 & .940 & .956 & 964 & .495 & .94 \\
\bottomrule
\end{tabular}}
\end{center}

Here Remain$_w$ is the progress-weighted global MAE and ``Remain primary'' is
the configured middle-progress weighted MAE; neither replaces the unweighted
remain MAE. The selected validation remain-MAE/cause-accuracy/cause-macro
triplets are $12.96/.931/.942$, $11.58/.928/.937$, and $11.64/.938/.939$.
\zh{表中的 Remain$_w$ 是按生产进度加权的全局 MAE，“Remain primary”是配置指定的中段进度加权 MAE；二者都不能替代未加权的剩余时间 MAE。Start$\leq5/10/15$ 入选 checkpoint 的验证集“剩余时间 MAE/原因准确率/原因宏召回率”分别为 $12.96/.931/.942$、$11.58/.928/.937$ 和 $11.64/.938/.939$。}

\begin{center}
\captionof{table}{No-clustering mixed-data per-class cause recall on test (four supported classes). / 无聚类混扰数据实验测试集逐类别原因召回率（四个有样本类别）。}
\scriptsize
\begin{tabular}{lrrrr}
\toprule
Start & Transport delay & Material shortage & Starved upstream & Queue buildup \\
\midrule
$\leq5$  & .962 & .989 & .906 & 1.000 \\
$\leq10$ & .928 & .956 & .893 & .984 \\
$\leq15$ & .943 & .978 & .910 & .995 \\
\bottomrule
\end{tabular}
\end{center}

Relative to the retained main model, mixed-data test F1 is lower at all three
horizons. The experiment therefore does not support the claim that directly
adding all available mixed-disturbance episodes improves the official score.
Because clustering is also disabled in the inherited recipe, the decline cannot
be attributed to either data mixing or clustering removal in isolation.
\zh{与最终保留的主模型相比，混扰数据实验在三个 Start 档位上的测试 F1 都更低。因此，现有结果不支持“直接加入全部混扰 episode 可以提高正式指标”这一结论。由于继承配方同时关闭了聚类，性能下降也不能单独归因于数据混合或去除聚类中的任一因素。}

\clearpage
\subsection{Graph and Station-Group Structure Ablations / 图结构与工位分组消融}
The \texttt{nograph} arm sets \texttt{graph\_mode=identity}, removing graph
edges. The \texttt{nogroup} arm removes station-group/semantic-group embedding.
\texttt{nopattern} disables the delayed pattern injection in the geographic
attention keys. All three arms use the same \texttt{dense\_i1} split as the final model but are
trained from scratch for at most 50 epochs with an empty \texttt{init\_ckpt}.
This budget differs from the warm-started 100-epoch full model and is reported
as a limitation of causal interpretation.
\zh{\texttt{nograph} 实验将 \texttt{graph\_mode} 设为 \texttt{identity}，从而移除图边；\texttt{nogroup} 实验移除工位分组/语义分组嵌入；\texttt{nopattern} 则关闭地理注意力键中的延迟 pattern 注入。三组消融都使用与最终模型相同的 \texttt{dense\_i1} 数据划分，但以空 \texttt{init\_ckpt} 从头训练，训练上限为 50 个 epoch。该预算不同于热启动且最多训练 100 个 epoch 的完整模型，因此这是解释消融因果效果时必须保留的限制。}

\begin{landscape}
\begin{center}
\captionof{table}{Complete validation/test event metrics for structure ablations; Report F1 is the strict station-plus-start metric. / 结构消融的完整验证与测试事件指标；Report F1 为“工位正确且起点满足容差”的严格指标。}
\scriptsize
\renewcommand{\arraystretch}{0.90}
\resizebox{0.97\linewidth}{!}{
\begin{tabular}{lllrrrrrrrrrrr}
\toprule
Arm & Start & Split & $P$ & $R$ & F1$_{\mathrm{who}}$ & F1$_{\mathrm{rep}}$ &
\uprep & \onwho & $n_{\rm up}$ & $n_{\rm on}$ & $n_{\rm all}$ & Best ep. \\
\midrule
Full    & $\leq5$  & Val  & .837 & .891 & .863 & .863 & .739 & .969 & \na & \na & \na & 39 \\
Full    & $\leq5$  & Test & .794 & .817 & .805 & .804 & .539 & .937 & 336 & 766 & 1102 & 39 \\
No graph& $\leq5$  & Val  & .823 & .658 & .731 & .731 & .115 & .937 & \na & \na & \na & 19 \\
No graph& $\leq5$  & Test & .832 & .678 & .747 & .746 & .101 & .930 & 336 & 766 & 1102 & 19 \\
No group& $\leq5$  & Val  & .782 & .765 & .774 & .774 & .383 & .962 & \na & \na & \na & 28 \\
No group& $\leq5$  & Test & .766 & .756 & .761 & .753 & .307 & .941 & 336 & 766 & 1102 & 28 \\
No pattern& $\leq5$  & Val  & .789 & .671 & .725 & .725 & .166 & .930 & \na & \na & \na & 15 \\
No pattern& $\leq5$  & Test & .799 & .698 & .748 & .745 & .158 & .935 & 336 & 766 & 1102 & 15 \\
\midrule
Full    & $\leq10$ & Val  & .831 & .840 & .835 & .835 & .711 & .960 & \na & \na & \na & 41 \\
Full    & $\leq10$ & Test & .806 & .761 & .783 & .778 & .528 & .936 & 604 & 768 & 1372 & 41 \\
No graph& $\leq10$ & Val  & .807 & .486 & .606 & .606 & .049 & .894 & \na & \na & \na & 9 \\
No graph& $\leq10$ & Test & .821 & .539 & .651 & .649 & .065 & .909 & 604 & 768 & 1372 & 9 \\
No group& $\leq10$ & Val  & .663 & .790 & .721 & .721 & .636 & .934 & \na & \na & \na & 44 \\
No group& $\leq10$ & Test & .605 & .745 & .668 & .651 & .445 & .947 & 604 & 768 & 1372 & 44 \\
No pattern& $\leq10$ & Val  & .598 & .659 & .643 & .627 & .360 & .939 & \na & \na & \na & 28 \\
No pattern& $\leq10$ & Test & .555 & .625 & .607 & .588 & .235 & .934 & 604 & 768 & 1372 & 28 \\
\midrule
Full    & $\leq15$ & Val  & .841 & .835 & .838 & .838 & .722 & .960 & \na & \na & \na & 42 \\
Full    & $\leq15$ & Test & .806 & .742 & .772 & .768 & .534 & .925 & 708 & 770 & 1478 & 42 \\
No graph& $\leq15$ & Val  & .819 & .435 & .568 & .568 & .028 & .883 & \na & \na & \na & 11 \\
No graph& $\leq15$ & Test & .851 & .482 & .616 & .612 & .037 & .887 & 708 & 770 & 1478 & 11 \\
No group& $\leq15$ & Val  & .620 & .760 & .683 & .683 & .609 & .927 & \na & \na & \na & 50 \\
No group& $\leq15$ & Test & .567 & .690 & .623 & .604 & .395 & .922 & 708 & 770 & 1478 & 50 \\
No pattern& $\leq15$ & Val  & .631 & .769 & .705 & .693 & .600 & .956 & \na & \na & \na & 42 \\
No pattern& $\leq15$ & Test & .539 & .681 & .616 & .601 & .397 & .947 & 708 & 770 & 1478 & 42 \\
\bottomrule
\end{tabular}}
\end{center}
\end{landscape}

\begin{center}
\captionof{table}{Complete remaining-time and cause metrics for the structure ablations. / 结构消融的完整剩余时间和原因识别指标。}
\scriptsize
\resizebox{\textwidth}{!}{
\begin{tabular}{lllrrrrrr}
\toprule
Arm & Start & Split & Remain MAE & Cause accuracy & Cause macro &
Majority baseline & Cause support & $\Delta$ test F1 vs full \\
\midrule
Full     & $\leq5$  & Val  & 12.3 & .956 & .963 & .409 & \na & \na \\
Full     & $\leq5$  & Test & 10.9 & .965 & .967 & .409 & 579 & .000 \\
No graph & $\leq5$  & Val  & 18.2 & .866 & .886 & .409 & \na & \na \\
No graph & $\leq5$  & Test & 17.0 & .876 & .889 & .409 & 579 & -.058 \\
No group & $\leq5$  & Val  & 13.2 & .932 & .952 & .409 & \na & \na \\
No group & $\leq5$  & Test & 12.7 & .943 & .955 & .409 & 579 & -.044 \\
No pattern & $\leq5$  & Val  & 20.1 & .726 & .725 & .409 & \na & \na \\
No pattern & $\leq5$  & Test & 18.6 & .732 & .734 & .409 & 579 & -.059 \\
\midrule
Full     & $\leq10$ & Val  & 12.6 & .958 & .968 & .409 & \na & \na \\
Full     & $\leq10$ & Test & 11.2 & .962 & .965 & .409 & 579 & .000 \\
No graph & $\leq10$ & Val  & 24.1 & .741 & .755 & .409 & \na & \na \\
No graph & $\leq10$ & Test & 23.3 & .734 & .746 & .409 & 579 & -.132 \\
No group & $\leq10$ & Val  & 12.5 & .928 & .944 & .409 & \na & \na \\
No group & $\leq10$ & Test & 11.6 & .938 & .953 & .409 & 579 & -.115 \\
No pattern & $\leq10$ & Val  & 14.8 & .807 & .823 & .409 & \na & \na \\
No pattern & $\leq10$ & Test & 13.4 & .793 & .805 & .409 & 579 & -.190 \\
\midrule
Full     & $\leq15$ & Val  & 11.0 & .967 & .978 & .409 & \na & \na \\
Full     & $\leq15$ & Test &  9.7 & .974 & .978 & .409 & 579 & .000 \\
No graph & $\leq15$ & Val  & 21.4 & .796 & .790 & .409 & \na & \na \\
No graph & $\leq15$ & Test & 20.8 & .794 & .787 & .409 & 579 & -.156 \\
No group & $\leq15$ & Val  & 11.1 & .941 & .950 & .409 & \na & \na \\
No group & $\leq15$ & Test & 10.6 & .938 & .949 & .409 & 579 & -.149 \\
No pattern & $\leq15$ & Val  & 12.5 & .851 & .862 & .409 & \na & \na \\
No pattern & $\leq15$ & Test & 11.1 & .843 & .850 & .409 & 579 & -.167 \\
\bottomrule
\end{tabular}}
\end{center}

\begin{landscape}
\begin{center}
\captionof{table}{Complete per-class cause recall for the full model and structure ablations on test. / 完整模型与结构消融在测试集上的逐类别原因召回率。}
\scriptsize
\begin{tabular}{llrrrrrr}
\toprule
Arm & Start & Transport delay & Material shortage & Starved upstream &
Queue buildup & Blocked downstream & High utilization \\
\midrule
Full     & $\leq5$  & .941 & .976 & .966 & .984 & no support & no support \\
No graph & $\leq5$  & .971 & 1.000 & .852 & .734 & no support & no support \\
No group & $\leq5$  & .956 & 1.000 & .911 & .952 & no support & no support \\
No pattern & $\leq5$  & .993 & 1.000 & .797 & .145 & no support & no support \\
\midrule
Full     & $\leq10$ & .926 & .988 & .962 & .984 & no support & no support \\
No graph & $\leq10$ & .993 & .963 & .738 & .290 & no support & no support \\
No group & $\leq10$ & .956 & 1.000 & .890 & .968 & no support & no support \\
No pattern & $\leq10$ & .941 & 1.000 & .802 & .476 & no support & no support \\
\midrule
Full     & $\leq15$ & .978 & .988 & .962 & .984 & no support & no support \\
No graph & $\leq15$ & .993 & .988 & .878 & .290 & no support & no support \\
No group & $\leq15$ & .926 & 1.000 & .911 & .960 & no support & no support \\
No pattern & $\leq15$ & .971 & 1.000 & .857 & .573 & no support & no support \\
\bottomrule
\end{tabular}
\end{center}
\end{landscape}

Removing graph structure causes the largest and most consistent degradation,
especially for upcoming recall, remaining-time MAE, and queue-buildup recall.
Removing group embeddings and delayed pattern injection also lowers F1 at
every horizon. The \texttt{nopattern} test report-F1 drops relative to the full
model are $.059/.190/.167$ for Start$\leq5/10/15$, respectively, and its
queue-buildup recall is substantially lower. However, the
\texttt{nogroup} validation precision values ($.782/.663/.620$) fail the
full-model $P\geq.80$ gate, and the precision gate also never opens for any
\texttt{nopattern} horizon; these rows therefore use the ungated validation
best. The \texttt{nogroup} Start15 checkpoint is the final epoch (50), so it
did not plateau within the assigned budget.
\zh{移除图结构造成最大且最一致的性能下降，尤其体现在 upcoming 召回率、剩余时间 MAE 和 queue-buildup 类召回率上。移除工位分组嵌入和延迟 pattern 注入也会降低三个 Start 档位的 F1。\texttt{nopattern} 在 Start$\leq5/10/15$ 下的测试 report-F1 相对完整模型分别下降 $.059/.190/.167$，且 queue-buildup 类召回率明显降低。不过，\texttt{nogroup} 的验证集精确率分别为 $.782/.663/.620$，均未满足完整模型的 $P\geq.80$ 门槛；\texttt{nopattern} 三档同样均未开启精确率门槛，因此这些行采用无门槛验证集最优值。\texttt{nogroup} 的 Start15 checkpoint 位于最后一个 epoch（50），说明该实验在给定预算内尚未进入平台期。}

The planned \texttt{nosem} arm
(\texttt{ablation\_nosem\_start5\_seed42}) aborted before epoch 1 because a
convolution weight had shape \texttt{[0,80,1,1]}; consequently, it has no valid
metrics. This failed arm is reported rather than silently excluded.
\zh{计划中的 \texttt{nosem} 实验（\texttt{ablation\_nosem\_start5\_seed42}）在第 1 个 epoch 之前即终止，原因是卷积权重形状为 \texttt{[0,80,1,1]}，因此没有有效指标。本文明确报告该失败实验，而不是将其静默删除。}

\clearpage
\section{Baseline Comparison / 基线对比}

\subsection{Comparable Setup and Limitations / 对齐设置与局限}
B4 is BTGCN with a GCN+GRU128 backbone; B5 is BSTAN with scalar
GAT+GRU128. Each baseline retains its own near-history input and parent
checkpoint and does not receive PHAST-DP-specific branches. Training follows
Start5$\rightarrow$Start10$\rightarrow$Start15. The six runs use the same
Start/Min8 semantics, 20 future grids, threshold candidates, 3-min start
tolerance, and validation checkpoint rule as the final PHAST-DP protocol.
All six use the fallback F1 rule because none meets both $P\geq.80$ and
$R\geq.70$.
\zh{B4 为采用 GCN+GRU128 骨干的 BTGCN，B5 为采用标量 GAT+GRU128 骨干的 BSTAN。每个基线保留自身的 near 历史输入和父 checkpoint，不引入 PHAST-DP 专有分支。训练顺序为 Start5$\rightarrow$Start10$\rightarrow$Start15。六组实验与最终 PHAST-DP 采用相同的 Start/Min8 语义、20 个未来网格、阈值候选集合、3 分钟起点容差以及验证集选模规则。由于六组都未同时满足 $P\geq.80$ 和 $R\geq.70$，最终均采用 F1 fallback 规则。}

The baseline split is 138/30/40 rather than 140/28/36, and the 40 baseline test
episodes have not been evaluated. Therefore, Tables~\ref{tab:b45-val} and
\ref{tab:b45-train} are complete validation/training reports, but they are not
test-set evidence that PHAST-DP outperforms B4/B5.
\zh{基线的数据划分为 138/30/40，而不是主模型的 140/28/36，并且基线的 40 个测试 episode 尚未参与评测。因此，表~\ref{tab:b45-val} 和表~\ref{tab:b45-train} 是完整的验证集/训练集报告，但不能作为 PHAST-DP 在测试集上优于 B4/B5 的证据。}

\begin{landscape}
\begin{center}
\captionof{table}{All official B4/B5 validation metrics; hit counts are over upcoming support. / B4/B5 全部正式验证指标；命中数均以 upcoming 支持数为分母。}
\label{tab:b45-val}
\scriptsize
\resizebox{0.97\linewidth}{!}{
\begin{tabular}{llrrrrllllrrr}
\toprule
Model & Start & $P$ & $R$ & F1 & \uprep & \upwho & \onwho &
Upcoming strict & Upcoming station & Best/total ep. & Threshold \\
\midrule
B4 & $\leq5$  & .8048 & .5655 & .6643 & .0742 & .0742 & .8078 &
32/431 & 32/431 & 1/41 & .50 \\
B5 & $\leq5$  & .8133 & .5709 & .6709 & .0812 & .0812 & .8124 &
35/431 & 35/431 & 2/42 & .70 \\
B4 & $\leq10$ & .8152 & .3829 & .5211 & .0353 & .0353 & .6833 &
27/764 & 27/764 & 3/43 & .90 \\
B5 & $\leq10$ & .8118 & .4163 & .5503 & .0314 & .0340 & .7466 &
24/764 & 26/764 & 2/42 & .85 \\
B4 & $\leq15$ & .6684 & .3718 & .4778 & .0280 & .0325 & .7133 &
25/892 & 29/892 & 5/45 & .94 \\
B5 & $\leq15$ & .8036 & .3543 & .4918 & .0291 & .0336 & .6772 &
26/892 & 30/892 & 2/42 & .90 \\
\bottomrule
\end{tabular}}
\end{center}
\end{landscape}

\begin{center}
\captionof{table}{All recorded B4/B5 training and auxiliary validation metrics; remain MAE is middle-progress-weighted in window units. / B4/B5 全部训练与辅助验证指标；剩余时间 MAE 为中段进度加权值，单位保留为窗口。}
\label{tab:b45-train}
\scriptsize
\resizebox{\textwidth}{!}{
\begin{tabular}{llrlrrrr}
\toprule
Model & Start & Train F1 & Train upcoming strict hits & Train \uprep &
Val remain MAE & Val cause accuracy & Val cause macro recall \\
\midrule
B4 & $\leq5$  & .6918 & 98/1753   & .0559 & 13.3421 & .8876 & .8631 \\
B5 & $\leq5$  & .7216 & 170/1753  & .0970 & 11.2884 & .9575 & .9538 \\
B4 & $\leq10$ & .6349 & 324/3109  & .1042 & 16.2363 & .9001 & .8750 \\
B5 & $\leq10$ & .6642 & 295/3109  & .0949 & 10.6643 & .9572 & .9481 \\
B4 & $\leq15$ & .8816 & 2657/3624 & .7332 & 11.7820 & .9040 & .8837 \\
B5 & $\leq15$ & .6434 & 477/3624  & .1316 &  9.6109 & .9518 & .9362 \\
\bottomrule
\end{tabular}}
\end{center}

B5 has higher validation F1 than B4 by $0.0066$, $0.0292$, and $0.0140$ at the
three horizons and lower remaining-time MAE in every pair. It does not dominate
strict upcoming recall: at Start$\leq10$, B5 obtains $0.0314$ versus B4's
$0.0353$. Both baselines remain weak on upcoming events. The most severe
generalization gap occurs for B4 Start$\leq15$, where train F1/upcoming strict
recall are $0.8816/0.7332$, but validation F1/upcoming strict recall fall to
$0.4778/0.0280$.
\zh{B5 在三个 Start 档位上的验证 F1 分别比 B4 高 $0.0066$、$0.0292$ 和 $0.0140$，且每一档的剩余时间 MAE 都更低。但 B5 并未在严格 upcoming 召回率上全面占优：Start$\leq10$ 时，B5 为 $0.0314$，低于 B4 的 $0.0353$。两个基线在 upcoming 事件上都较弱。最严重的泛化差距出现在 B4 Start$\leq15$：训练 F1/严格 upcoming 召回率为 $0.8816/0.7332$，验证集则降至 $0.4778/0.0280$。}

\subsection{Overall Discussion and Validity Boundaries / 总体讨论与有效性边界}
The retained \texttt{12\_3} PHAST-DP reaches test station-level F1 of
$0.805/0.783/0.772$ and strict upcoming recall of
$0.539/0.528/0.534$ for Start$\leq5/10/15$. The graph and station-group
ablations both reduce F1, and graph removal nearly eliminates upcoming recall at
long horizons. The no-clustering mixed-data experiment also underperforms the
retained model, but it changes both data composition and clustering, so it
cannot isolate either factor.
\zh{最终保留的 \texttt{12\_3} PHAST-DP 在 Start$\leq5/10/15$ 三档上的测试集工位级 F1 分别为 $0.805/0.783/0.772$，严格 upcoming 召回率分别为 $0.539/0.528/0.534$。移除图结构或工位分组都会降低 F1，其中移除图结构几乎使长预测范围的 upcoming 召回能力消失。无聚类混扰数据实验也弱于最终模型，但由于它同时改变了数据组成和聚类设置，因此无法单独识别任何一个因素的影响。}

Four validity limits must accompany these results. First, all runs use one seed,
so uncertainty and significance are unknown. Second, B4/B5 currently have no
test result and use a different episode split; their validation scores cannot be
presented as a final test ranking. Third, the structure ablations use a smaller,
from-scratch budget than the warm-started full model. Fourth, remaining-time
metrics are not uniformly comparable: PHAST-DP reports minute-scale
\texttt{remain\_len\_mae}, while B4/B5 retain middle-progress-weighted values in
window units. These boundaries are preserved in the tables rather than being
hidden by a single pooled score.
\zh{这些结果必须同时说明四项有效性限制。第一，所有实验只使用一个随机种子，因此无法估计不确定性或统计显著性。第二，B4/B5 尚无测试集结果，且采用不同的 episode 划分，所以不能将其验证分数作为最终测试排名。第三，结构消融采用比完整模型更小的、从头训练的预算。第四，剩余时间指标并非完全可比：PHAST-DP 报告以分钟为单位的 \texttt{remain\_len\_mae}，B4/B5 则保留以窗口为单位的中段进度加权值。本文在各表中明确保留这些边界，而不是用一个混合总分掩盖差异。}

\end{document}
